% --- 1. INFORMAZIONI GENERALI ---
\section{Informazioni Generali}
\begin{itemize}
	\item \textbf{Luogo/Piattaforma:} {Server Discord}
	\item \textbf{Data:} 2026-03-20
	\item \textbf{Orario di inizio:} 15:00
	\item \textbf{Orario di fine:} 17:00
	\item \textbf{Responsabile:} {Sofia De Blasi}
	\item \textbf{Scriba:} {Francesco Marcon}
\end{itemize}

\vspace{0.5cm}
\textbf{Partecipanti presenti:}
\begin{table}[ht]
	\centering
	\renewcommand{\arraystretch}{1.2}
	\begin{tabularx}{\textwidth}{X l}
		\toprule
		\textbf{Nome e Cognome} & \textbf{Matricola} \\
		\midrule
		Sofia De Blasi & 2111014 \\
		Roberto Mariano Doroftei & 2111031 \\
		Filippo Idiometri & 2035617 \\
		Francesco Marcon & 2110990 \\
		Armando Moda Scarati & 2082864 \\
		Anton Ricardo Rupi & 2054304 \\
		\bottomrule
	\end{tabularx}
\end{table}

\vspace{0.5cm}
\textbf{Partecipanti assenti:}
\begin{table}[ht]
	\centering
	\renewcommand{\arraystretch}{1.2}
	\begin{tabularx}{\textwidth}{X l}
		\toprule
		\textbf{Nome e Cognome} & \textbf{Matricola} \\
		\midrule
		Nessuno & - \\
		\bottomrule
	\end{tabularx}
\end{table}

% --- 2. ORDINE DEL GIORNO ---
\section{Ordine del Giorno}
Gli argomenti previsti per la riunione odierna sono:
\begin{enumerate}
	\item Revisione e approvazione verbale esterno per comunicazione mail con l'azienda Zucchetti, invio mail con richiesta di firma e proposta di riunione da remoto
	\item Comunicazione via mail con azienda Var Group.
	\item Riorganizzazione delle norme di progetto secondo lo standard ISO/IEC 12207:1995
	\item Standard interni di gestione e nomenclatura di branch, task e versionamento.
	\item Aggiornamento sito e documentazione.
	\item Ordinamento documenti.
	\item Protezione del branch main sulla repository GitHub.
	\item Revisione del preventivo costi e delle ore di progetto.
	\item Allineamento di strumenti di sviluppo interni, e workflow Latex.
\end{enumerate}

% --- 3. RESOCONTO E DISCUSSIONE ---
\section{Resoconto della Riunione}
Durante la riunione il gruppo ha discusso principalmente aspetti legati alla gestione della documentazione, alla comunicazione con le aziende proponenti (in particolare Zucchetti e VarGroup), all’organizzazione della repository e del sito, oltre che alla definizione di alcune convenzioni operative e alla revisione dei costi di progetto.

Sono state inoltre definite diverse convenzioni operative su versionamento, gestione dei branch, aggiornamento del sito e struttura delle Norme di Progetto, con l’obiettivo di allinearsi agli standard richiesti e ridurre il rischio di errori in fase di valutazione.

Infine, il gruppo ha rivisto la distribuzione delle ore e dei costi del progetto, trovando un compromesso più adeguato.

\subsection{Organizzazione Riunioni con Zucchetti}
Il gruppo ha ritenuto necessario proporre all'azienda Zucchetti una riunione conoscitiva da remoto, per chiarire alcuni aspetti tecnici che la comunicazione via mail non ha sufficientemente chiarito.

\textbf{Motivazione:} ricevere ulteriori specifiche tecniche sul progetto e conoscere di persona i referenti dell'azienda Zucchetti.

Azioni operative:
\begin{itemize}
	\item Invio di un'ulteriore mail a Zucchetti con:
	\begin{itemize}
		\item Richiesta di firma del verbale precedente.
		\item Proposta di riunione (da remoto) entro la settimana.
		\item Includere una serie di domande preparate dal gruppo.
	\end{itemize}
\end{itemize}

\subsection{Comunicazione con Var Group}
È stato evidenziato che VarGroup risponde lentamente alle mail o non risponde affatto.

\textbf{Motivazione:} necessità di ottenere risposte per valutare più approfonditamente anche la fattibilità del secondo capitolato scelto dal gruppo.

Azioni operative:
\begin{itemize}
	\item Reinvio della mail già preparata per sollecitare una risposta.
\end{itemize}


\subsection{Riorganizzazione Norme di Progetto Secondo Standard ISO}
Il gruppo ha discusso la struttura delle Norme di Progetto, attualmente non allineata allo standard ISO.

\textbf{Motivazione:} allineamento effettivo secondo lo standard definito dall'ISO/IEC 12207:1995.

Azioni operative:
\begin{itemize}
	\item Riorganizzazione della struttura secondo le categorie previste dallo standard ISO, mantenendo i contenuti già prodotti ma redistribuendoli in macro-sezioni corrette.
\end{itemize}

\subsection{Gestione Branch e Versionamento}
È stata affrontata la gestione dei branch nella repository GitHub.

\textbf{Motivazione:} garantire un flusso di lavoro ordinato e coerente.

Azioni Operative e Convenzioni Stabilite:
\begin{itemize}
	\item Eliminazione dei branch personali.
	\item Creazione di un branch per ogni task.
	\item Utilizzo della seguente convenzione di nomenclatura:
	\begin{itemize}
		\item \textbf{t-xy}: identificativo della task.
		\item \textbf{gestionale} o \textbf{tecnico}: tipo del documento.
		\item \textbf{categoria}: categoria della task (esempio: aggiornamento, creazione, rimozione).
		\item \textbf{nomeDocumento}: nome del documento trattato.
	\end{itemize}
\end{itemize}

\subsection{Aggiornamento Sito e Documentazione}
Il gruppo ha definito le modalità di aggiornamento del sito.

\textbf{Motivazione:} evitare incoerenze tra documentazione e sito web.

Convenzioni stabilite:
\begin{itemize}
	\item Aggiornamento obbligatorio di versione e data di modifica sul sito a ogni modifica di documento
	\item Verifica dell’ordinamento corretto dei documenti
\end{itemize}

\subsection{Ordinamento Documenti nel Sito}
Sono state definite le regole di ordinamento nel sito.

\textbf{Motivazione:} garantire coerenza e facilità di consultazione.

Convenzioni stabilite:
\begin{itemize}
	\item Verbali ordinati in LIFO per data della riunione.
	\item Altri documenti ordinati in LIFO per data di ultima modifica.
\end{itemize}

\subsection{Protezione del Branch Main}
È stata discussa la gestione del branch principale sulla repository GitHub.

\textbf{Motivazione:} garantire controllo e qualità del codice/documentazione.

Convenzioni stabilite:
\begin{itemize}
	\item Divieto di push diretto su main.
	\item Aggiornamento solo tramite pull request.
	\item Obbligo di almeno una revisione approvata.
\end{itemize}

\subsection{Revisione Costi e Ore di Progetto}
Il gruppo ha rivisto la distribuzione delle ore e dei costi.

\textbf{Motivazione:} il costo risultava troppo elevato rispetto ai riferimenti.

Azioni operative:
\begin{itemize}
	\item Ridefinizione delle ore per ruolo per ottenere un compromesso tra costo e distribuzione equilibrata del lavoro.
\end{itemize}

\subsection{Allineamento strumenti e workflow Latex}
È stato fatto un allineamento sugli strumenti per la gestione dei documenti.

\textbf{Motivazione:} uniformare il workflow e semplificare la collaborazione.

Azioni operative:
\begin{itemize}
	\item Utilizzo di Latex con compilazione automatica (Latex Workshop) e condivisione di linee guida comuni.
\end{itemize}

% --- 4. DECISIONI PRESE ---
\section{Decisioni Prese}
\renewcommand{\arraystretch}{1.3}
\vspace{-0.3cm}
\noindent
\begin{longtable}{c p{12cm}}
	\toprule
	\textbf{Codice} & \textbf{Decisione} \\
	\midrule
	\endfirsthead

	\toprule
	\textbf{Codice} & \textbf{Decisione} \\
	\midrule
	\endhead

	\midrule
	\multicolumn{2}{r}{\textit{Continua nella pagina successiva}} \\
	\midrule
	\endfoot

	\bottomrule
	\endlastfoot

	\textbf{VI-4.1} & Inviare mail a Zucchetti per firma verbale e richiesta nuova riunione. \\
	\textbf{VI-4.2} & Reinvio mail a VarGroup per sollecito risposta. \\
	\textbf{VI-4.3} & Riorganizzare le Norme di Progetto secondo standard ISO. \\
	\textbf{VI-4.4} & Eliminare branch personali e usare branch per task. \\
	\textbf{VI-4.5} & Aggiornare sempre sito (versione e data) dopo modifiche ai documenti. \\
	\textbf{VI-4.6} & Definire ordinamento LIFO per documenti e verbali. \\
	\textbf{VI-4.7} & Proteggere branch main con pull request e revisione obbligatoria. \\
	\textbf{VI-4.8} & Ricalibrare ore e costi del progetto. \\
	\textbf{VI-4.9} & Standardizzare utilizzo strumenti Latex. \\
	\textbf{VI-4.10} & Impostare date progetto: inizio 2026-03-27, fine 2026-08-28. \\
\end{longtable}

% --- 5. AZIONI (TODO) ---
\section{Azioni da Svolgere (To-Do)}
\begin{longtable}{c c p{4.5cm} l l}
	\toprule
	\textbf{Codice} & \textbf{Rif.} & \textbf{Descrizione Task} & \textbf{Assegnatario} & \textbf{Scadenza} \\
	\midrule
	\endfirsthead

	\toprule
	\textbf{Codice} & \textbf{Rif.} & \textbf{Descrizione Task} & \textbf{Assegnatario} & \textbf{Scadenza} \\
	\midrule
	\endhead

	\midrule
	\multicolumn{5}{r}{\textit{Continua nella pagina successiva}} \\
	\midrule
	\endfoot

	\bottomrule
	\endlastfoot
	
	\textbf{T-4.1} & VI-4.1 & Invio mail a Zucchetti (firma verbale + richiesta riunione) & Francesco & 2026-03-23 \\
	\textbf{T-4.2} & VI-4.1 & Preparazione e divisione domande per riunione & Sofia, Ricardo, Francesco & 2026-03-24 \\
	\textbf{T-4.3} & VI.4.2 & Reinvio mail a VarGroup & Ricardo & 2026-03-23 \\
	\textbf{T-4.4} & VI-4.3 & Riorganizzazione Norme di Progetto & Sofia & 2026-03-25 \\
	\textbf{T-4.5} & VI.4.3 & Revisione Norme di Progetto & Filippo & 2026-03-26 \\
	\textbf{T-4.6} & VI.4.5 & Aggiornamento sito per documenti modificati & Tutti & - \\
	\textbf{T-4.7} & VI.4.4 & Applicazione nuova gestione branch & Tutti & - \\
	\textbf{T-4.8} & VI.4.7 & Applicazione revisione obbligatoria Pull Request & Tutti & - \\
	\textbf{T-4.9} & VI.4.8 & Aggiornamento documento dichiarazione impegni & Armando & 2026-03-25 \\
	\textbf{T-4.10} & VI.4.8 & Revisione dichiarazione impegni & Sofia & 2026-03-26 \\
	\textbf{T-4.11} & VI.4.5 & Pubblicazione valutazione capitolati & Filippo & 2026-03-25 \\
	\textbf{T-4.12} & VI.4.5 & Revisione valutazione capitolati & Armando & 2026-03-26 \\
	\textbf{T-4.13} & VI.4.5 & Redazione verbale esterno VarGroup (parziale, in attesa di risposta) & Ricardo & 2026-03-25 \\
	\textbf{T-4.14} & VI.4.8 & Redazione lettera di presentazione & Roberto & 2026-03-25 \\
\end{longtable}

% --- 6. PROSSIMA RIUNIONE ---
\section{Prossima Riunione}
\begin{itemize}
	\item \textbf{Data:} {2026-03-24}
	\item \textbf{Ora:} {14:30}
	\item \textbf{OdG preliminare:} Revisione generale dei documenti per la candidatura, verifica di eventuale risposta da parte delle aziende Zucchetti e Var Group.
\end{itemize}